\infographicchapter{Corneal Topography}{corneal_topography.png}
\label{ch:infographic-corneal-topography}

\section{Overview}

Corneal topography (also called photokeratoscopy or videokeratography) is a quick, completely non-touch eye test that creates a detailed color-coded map of the front curve of your cornea (the clear dome at the front of your eye). Because the cornea provides roughly 70\% of your eye's focusing power, even small shape changes can blur vision; much like a geographical map charts hills and valleys, corneal topography reveals subtle steep and flat zones that a standard eye exam cannot see. Most standalone topographers use a clever technique called a Placido disc: a lighted cone projects a set of glowing concentric rings onto your eye. Because your smooth tear film acts like a tiny curved mirror, a digital camera snaps a picture of how those rings reflect back. Where the cornea is steeper, the reflected rings sit closer together; where it is flatter, they spread farther apart. A computer turns that ring pattern into an easy-to-read color map of your cornea's focusing power.

Eye doctors use corneal topography to measure astigmatism precisely, spot and monitor corneal bulging conditions such as keratoconus and pellucid marginal degeneration, evaluate eyes before and after laser vision correction, and fit custom contact lenses (including rigid gas-permeable, scleral, and overnight reshaping lenses). It also helps map uneven corneal curves caused by scars, sun-related surface growths (pterygia), corneal transplants, or past eye injuries, and even checks how smoothly your tear film holds together in dry eye disease.

Its biggest strengths are speed, zero eye contact, high sensitivity to subtle front-surface bumps or warps, easy office use, and complete safety with no X-rays. Its main limit is that standard Placido topographers only see the front surface of the cornea---they cannot measure the back surface or corneal thickness, which are needed to catch the very earliest stages of keratoconus and require corneal tomography instead. Results can also be thrown off by dry eye, surface scratches, looking slightly off-center, or shadows from brows and lashes.

\section{Underlying Physical Principles}
\begin{itemize}
    \item \textbf{Projecting Glowing Rings (Placido Disc):} A small lighted bowl or cone shines a neat pattern of alternating bright and dark circles (called mires) onto the front of your eye using safe visible or near-infrared LEDs.
    \item \textbf{Your Tear Film Acts Like a Mirror:} The smooth surface where air meets your thin layer of tears reflects light just like a curved silver mirror, sending a crisp reflection of the circles straight back toward a camera in the center of the cone.
    \item \textbf{How Ring Spacing Reveals Shape:} A steeper curve pulls neighboring reflected rings tightly together (stronger focusing power), while a flatter curve spreads them farther apart (weaker focusing power). Irregular shapes bend the circles into ovals or waves.
    \item \textbf{Measuring Thousands of Points:} A central digital camera photographs the reflected rings, and software traces the edges of every circle along thousands of spokes to calculate the exact slope and curvature at each spot.
    \item \textbf{Turning Curves Into Color Maps:} The computer converts those measurements into focusing power (diopters) and paints a color map (most commonly a sagittal curvature map)---typically using warm reds and oranges for steeper areas and cool blues and greens for flatter zones---while also calculating simulated keratometry (SimK) values that mirror traditional keratometer readings.
\end{itemize}

\section{Imaging Mechanism}
\begin{enumerate}
    \item \textbf{Lining Up Without Touching:} You rest your chin and forehead on the supports while the operator (or an automatic tracker) centers the lighted ring cone at the exact right distance in front of your eye.
    \item \textbf{Fresh Blink and Quick Snapshot:} You blink once or twice to smooth your tear film, then look straight ahead as the glowing rings reflect off your eye during a 1--2 second snapshot.
    \item \textbf{Digital Ring Capture:} A camera sitting at the center of the cone captures the reflected circle pattern in sharp detail and checks that your eye is centered and in focus.
    \item \textbf{Computer Ring Tracing:} Software locates the edge of each ring along hundreds of radial lines, compares the spacing to an ideal smooth curve, and calculates the focusing power at thousands of points.
    \item \textbf{Color Map Display:} The workstation displays color-coded curvature maps along with automatic astigmatism and keratoconus scores for your eye doctor to review right away.
\end{enumerate}

\section{Key Technical Parameters}
\noindent
\begin{tabularx}{\textwidth}{@{}l X@{}}
\toprule
\textbf{Parameter} & \textbf{Typical Value / Description} \\
\midrule
Physical Principle & Glowing concentric rings (Placido disc) reflect off the eye's tear film to map front curvature \\
Energy / Signal Source & Visible or near-infrared LED ring projector and digital camera \\
Typical Frequency / Energy & Visible white or red/near-infrared light (about 630--950 nm; 16--32 concentric rings) \\
Spatial Resolution & About 0.1--0.25 diopter focusing detail across 7,000--10,000+ points on the cornea \\
Typical Scan Duration & 1--2 seconds per eye (under 5 minutes total for both eyes) \\
Contrast Agent & None (completely non-touch light reflection) \\
Ionizing Radiation & None \\
\bottomrule
\end{tabularx}

\section{Common Symptoms That Prompt This Scan}
\begin{itemize}
    \item \textbf{Frequently Changing Glasses Prescriptions or Worsening Blur:} Increasing astigmatism, ghosting around letters, or glare and halos at night that may signal a bulging cornea (keratoconus).
    \item \textbf{Uncomfortable Contact Lens Fit:} Lenses slipping, irritating the eye, or failing to give crisp vision due to an irregular front eye surface.
    \item \textbf{Planning Laser Vision Correction or Cataract Surgery:} Routine pre-surgical surface mapping---even with no eye disease---to plan LASIK, PRK, or specialty astigmatism-correcting lenses.
    \item \textbf{Blurry Vision After Eye Injury or Corneal Surgery:} Evaluating surface scars, uneven healing, or tight stitches affecting the front curve of the eye.
\end{itemize}

\section{Clinical Usages}

\subsection{Primary Usages}
\begin{itemize}
    \item \textbf{Measuring Astigmatism:} Maps the exact strength, angle, and symmetry of corneal astigmatism, telling ordinary regular astigmatism apart from uneven warping or eye disease.
    \item \textbf{Spotting and Tracking Keratoconus:} Detects classic front-surface warning patterns---such as uneven lower steepening, angled steep axes, or a ``crab-claw'' shape---in keratoconus and pellucid marginal degeneration, and compares repeat scans over time to see whether bulging is stable or actively progressing (helping guide stabilizing treatments like corneal cross-linking).
    \item \textbf{Laser Eye Surgery Checks:} Checks front-surface symmetry before laser vision correction and evaluates how smoothly and centrally the treated zone healed afterward.
    \item \textbf{Fitting Specialty Contact Lenses:} Guides custom lens selection for rigid gas-permeable, scleral, and overnight orthokeratology lenses, and detects temporary corneal warping caused by tight or worn-out lenses.
    \item \textbf{Mapping Uneven Corneas:} Pinpoints surface distortion caused by scars, corneal swelling, fleshy growths (pterygia), raised nodules, or past eye injuries.
\end{itemize}

\subsection{Secondary Usages}
\begin{itemize}
    \item \textbf{Guiding Stitch Removal After Corneal Transplant:} Finds the steepest pulling axis after a corneal transplant so the surgeon knows which tiny stitches to remove to smooth out astigmatism.
    \item \textbf{Dry Eye and Tear-Film Checks:} Measures how quickly the smooth ring reflection breaks up between blinks, giving a gentle, non-touch check of tear-film stability in dry eye disease.
    \item \textbf{Cataract Astigmatism Alignment:} Confirms the front astigmatism angle when planning relaxing incisions or toric cataract lens implants, usually paired with corneal tomography.
    \item \textbf{Investigating Glare and Halos:} Estimates subtle front-surface focusing imperfections (higher-order aberrations) when patients notice unexplained glare, halos, or ghosting.
\end{itemize}

\section{Advantages}
\begin{itemize}
    \item \textbf{Highly Sensitive to Front-Surface Changes:} Quickly catches tiny front-curve steepening, uneven astigmatism, and contact lens warping on intuitive color maps.
    \item \textbf{Fast, Non-Touch, and Completely Safe:} Captures pictures in a couple of seconds using ordinary reflected light---no eye contact, no drops, and no X-rays.
    \item \textbf{Widely Available and Cost-Effective:} Standalone topographers cost less than full 3D tomographers and are found in many optometry and ophthalmology offices.
    \item \textbf{Detailed Astigmatism and Optical Metrics:} Provides exact axis, symmetry, and surface-regularity measurements helpful for lens fitting and post-transplant stitch care.
    \item \textbf{Bonus Tear-Film Check:} Watching how crisp the rings stay between blinks lets doctors evaluate dry eye without touching the eye or adding dye.
\end{itemize}

\section{Disadvantages}
\begin{itemize}
    \item \textbf{Sees Only the Front Surface:} Standard Placido topographers cannot measure the back surface of the cornea or corneal thickness, so they can miss very early keratoconus unless paired with corneal tomography.
    \item \textbf{Depends Heavily on Smooth Tears:} Because rings bounce off your tear film, dry spots, surface scratches, or watery eyes immediately break up the rings and create false irregular patterns.
    \item \textbf{Tiny Central Gap and Edge Shadows:} The small spot directly behind the camera lens must be estimated by software, while deep-set eyes, prominent noses, lashes, or steep cones can cast shadows near the outer edges.
    \item \textbf{Affected by Aim and Contact Lenses:} Looking slightly off-center can create false lower steepening (``pseudokeratoconus''), and recent contact lens wear can temporarily mold the cornea into misleading shapes.
    \item \textbf{Struggles With Heavy Scars:} Badly scarred, ulcerated, or very rough corneas scatter the reflected rings, making accurate mapping difficult or impossible.
\end{itemize}

\section{Contraindications}

\subsection{Absolute}
\begin{itemize}
    \item \textbf{None:} Because Placido corneal topography never touches the eye, uses gentle visible or near-infrared light, and involves no radiation or dye, it has no absolute medical contraindications.
\end{itemize}

\subsection{Relative}
\begin{itemize}
    \item \textbf{Contagious Eye Infection:} Non-urgent scans should be delayed during contagious pink eye or corneal infection (such as adenoviral infection) to prevent spreading germs on clinic equipment.
    \item \textbf{Recent Contact Lens Wear or Severe Surface Damage:} Scanning without a proper break from contact lenses can produce misleading warpage, and severe corneal scarring or surface breakdown may prevent clear ring reflection.
\end{itemize}

\section{Patient Safety Guidelines}
\begin{itemize}
    \item \textbf{Safe, Non-Touch Light:} Placido rings stay well within FDA and ANSI eye light-safety limits with zero X-ray radiation; staff keep the lighted cone at a safe distance so it never bumps the eye.
    \item \textbf{Equipment Disinfection:} Following AAO hygiene standards, staff clean chinrests, forehead bands, and handles between every patient and reschedule elective scans during active contagious eye infections.
    \item \textbf{Never Use Alone for Laser Surgery Clearance:} Because Placido topography cannot see the back of the cornea, doctors always pair it with 3D corneal tomography before clearing anyone for laser vision surgery and enforce an adequate contact lens break.
    \item \textbf{Gentle Eyelid and Tear Care:} Technicians encourage natural blinking between snapshots and avoid pressing on the eyeball if lifting a droopy upper lid, protecting fragile surface cells and keeping the map accurate.
    \item \textbf{Checking Raw Ring Quality:} Clinicians inspect the actual reflected rings before making treatment choices so dry spots or off-center gaze are never mistaken for keratoconus.
\end{itemize}

\section{Prerequisite Tests}
\begin{itemize}
    \item \textbf{No Blood Tests Needed:} Because this is a non-touch light scan, no lab work, kidney tests, or special medical clearances are required.
    \item \textbf{Standard Eye Exam First:} Per AAO guidelines, your eye doctor first checks your vision, glasses prescription, and front-of-eye health under a slit-lamp microscope to spot scars, surface growths, or dry eye that could warp the rings.
    \item \textbf{Contact Lens Break Verification:} Before laser-surgery screening or keratoconus testing, staff confirm you have stayed out of contact lenses long enough (usually 1--2 weeks for soft lenses and 2--4 weeks for rigid lenses) so temporary molding fades.
    \item \textbf{Pairing With Tomography for Surgery:} When screening for laser vision correction or early keratoconus, corneal tomography is ordered alongside topography to check thickness and the back surface.
\end{itemize}

\section{Preparing for the Scan}
\begin{itemize}
    \item \textbf{Follow Your Contact Lens Break:} Take out contact lenses before the visit. For laser surgery screening or keratoconus evaluation, strictly follow your doctor's lens-free waiting period (typically 1--2 weeks for soft lenses and 2--4 weeks for rigid lenses) so your cornea returns to its natural shape.
    \item \textbf{No Fasting or Eye Drops Needed:} Eat, drink, and take regular medications normally; no drops are needed for the scan itself.
    \item \textbf{Keep Eyelids Clean and Skip Thick Gels:} Avoid heavy eye makeup, mascara, or oily cleansers around your lids, and do not put in thick lubricating gels or ointments right before the scan, since greasy tears blur the reflected rings.
    \item \textbf{Blink and Look Straight Ahead:} You will rest your chin and forehead on the clean supports, blink once or twice right before the photo to smooth your tears, then open wide and look straight at the center target light so the map stays perfectly centered.
    \item \textbf{No Driver Needed:} Standalone corneal topography does not blur vision or use sedation, so you can drive home right away (unless dilating drops are used for another part of your visit).
\end{itemize}

\section{During Scan}
\begin{itemize}
    \item \textbf{Chinrest and Center Target:} You sit upright with your chin on the rest and forehead lightly against the strap, looking straight ahead at the glowing target in the middle of the lighted cone.
    \item \textbf{Easy Breathing and Fresh Blink:} Breathe quietly and naturally. Blink fully right before each snapshot when asked, then hold your eyes wide open without squinting for the quick picture.
    \item \textbf{What You Will See and Feel:} Nothing touches your eye, and the test is silent and painless. You simply see bright concentric circles reflected in front of you with no pressure and no drops.
    \item \textbf{Let Staff Know If Uncomfortable:} Tell the technician if your eyes feel dry, watery, or tired from holding them open, or if the chinrest height feels awkward, so they can pause or adjust the chair.
    \item \textbf{How Long It Takes:} Each snapshot takes only 1--2 seconds, and checking both eyes---including a couple of extra shots to make sure the rings are crisp---usually takes less than 5 minutes.
\end{itemize}

\section{Normal Appearance}
\begin{itemize}
    \item \textbf{Crisp Reflected Rings:} Raw photos show neat, unbroken, evenly spaced concentric circles reflected cleanly across the front of the cornea without wavy spots or dry-eye gaps.
    \item \textbf{Balanced Color Shapes:} The color curvature map shows a smooth round, oval, or symmetrical bowtie pattern (vertical or horizontal for ordinary astigmatism) in calm green and yellow shades, with normal focusing power around 41--46 diopters.
    \item \textbf{Top-to-Bottom Balance:} The upper and lower halves of the cornea match closely (usually within 1.4 diopters of each other), and the steep and flat directions cross cleanly at right angles without crooked axes.
\end{itemize}

\section{Observing for Abnormal Patterns}
\begin{itemize}
    \item \textbf{Uneven Bowtie With Crooked Axes (Keratoconus):} A vertical bowtie that is noticeably steeper on the bottom half than the top---especially when the upper and lower halves tilt at an angle over 21 degrees from each other---is a classic front-surface sign of keratoconus.
    \item \textbf{``Crab-Claw'' Pattern:} Strong steepening curving along the lower outer edges (around the 4- and 8-o'clock positions) with a flatter strip in the middle creates a ``crab-claw'' or ``kissing-doves'' look classically seen in pellucid marginal degeneration (though lower keratoconus can look similar and needs a thickness map to confirm).
    \item \textbf{Patterns After Laser Eye Surgery:} Past nearsighted (myopic) laser surgery leaves a neat, centered blue flattened circle surrounded by steeper outer tissue, while farsighted surgery leaves a steeper central island. An off-center zone shows decentered treatment, and a new steep spot inside a flattened zone can signal post-LASIK bulging.
    \item \textbf{Localized Uneven Spots:} Abrupt, lopsided flat or steep patches point to surface issues such as a fleshy growth on the white of the eye (pterygium), raised nodules, corneal scars, or tight transplant stitches.
    \item \textbf{Harmless Look-Alikes:} Looking slightly off-center, a heavy upper eyelid, or recent contact lens molding can mimic keratoconus (``pseudokeratoconus''), while dry eye, mucus threads, or surface scratches break up the rings and mimic irregular astigmatism.
\end{itemize}

\section{Dealing with Uncertainties}
\begin{itemize}
    \item \textbf{Dry Eye and Missing Data:} Dry spots, surface scratches, watery eyes, or shadows from brows and lashes immediately disrupt the reflected rings and leave blank patches or false irregular shapes.
    \item \textbf{False Keratoconus Look-Alikes:} Off-center gaze, eyelid pressure, and recent contact lens wear often produce uneven lower steepening that looks remarkably like true keratoconus.
    \item \textbf{Cannot See the Back Surface:} Because Placido topographers only measure the front surface and estimate the tiny spot behind the camera, a normal front map cannot rule out early back-surface bulging or corneal thinning.
    \item \textbf{Checking Raw Rings and Retaking:} Doctors sort out doubtful maps by inspecting the raw circle photos for blur or off-center aim, calming dry eye with drops, extending the contact lens break, and taking fresh centered snapshots to confirm consistency.
    \item \textbf{Adding Tomography and Follow-Up:} Any borderline uneven bowtie, crooked axes, or laser-surgery screening is paired with 3D corneal tomography. If contact lens molding is suspected, mapping is repeated after 2--4 weeks lens-free to make sure the cornea returns to normal.
\end{itemize}

\section{Augmenting Methods}
\begin{itemize}
    \item \textbf{3D Corneal Tomography (Scheimpflug or AS-OCT):} Adds back-surface elevation and full-cornea thickness maps, filling the exact blind spot of front-only Placido topography for early keratoconus and laser-surgery safety.
    \item \textbf{Wavefront Aberrometry:} Measures focusing imperfections across the entire eye (not just the cornea) to tell whether glare or halos come from the cornea, the natural lens, or deeper inside.
    \item \textbf{Dry Eye Testing:} Tear-saltiness checks, dye tear-breakup exams, and eyelid oil-gland imaging uncover dry eye conditions that distort reflected rings and guide treatment before repeat scans.
    \item \textbf{Optical Biometry:} Measures eye length and total corneal power alongside topography to sharpen cataract lens-implant calculations.
    \item \textbf{Automated Keratoconus Scores and Lens Simulators:} Built-in scoring tools (like KISA\% and top-to-bottom difference scores) standardize ectasia screening, while lens-fitting software turns the color map directly into custom rigid or scleral lens designs.
\end{itemize}

\section{Comparison with Alternative Modalities}
\begin{itemize}
    \item \textbf{Versus Corneal Tomography:} Placido topography is faster, less expensive, and super sensitive to tiny front-surface and tear-film ripples. Corneal tomography adds back-surface shape and full thickness, making it mandatory for laser-surgery screening and early keratoconus detection.
    \item \textbf{Versus Basic Keratometry:} Standard keratometry measures only a small 3 mm ring near the center and assumes a smooth eye, while topography maps thousands of points across 8--10 mm to catch irregular curves and off-center cones.
    \item \textbf{Versus Anterior Segment OCT:} AS-OCT sees through scars and directly measures individual corneal layers, whereas Placido topography depends on mirror-like tear reflection and struggles on heavily scarred surfaces.
\end{itemize}

\section{Follow-up Steps}
\begin{itemize}
    \item \textbf{Right After the Scan:} No recovery time is needed---you can return to normal activities right away (staying out of contact lenses only if your clinic asked you to wait for repeat scans or surgical evaluation).
    \item \textbf{Double-Checking Map Quality:} Your clinician checks the raw ring photos for dry spots, lid shadows, or off-center gaze and pairs the map with your eye exam and prescription. If dry eye blurred the scan, treating dryness and repeating the map gives a cleaner answer.
    \item \textbf{Confirming With 3D Tomography:} Any uneven lower steepening, crooked axes, crab-claw shape, or laser-surgery screening is followed up with 3D corneal tomography to check thickness and the back surface.
    \item \textbf{Targeted Follow-Up Visits:} Suspected contact lens warping is rechecked after 2--4 weeks lens-free; post-transplant maps guide selective stitch removal; and stable keratoconus or custom contact lens fits are typically rechecked every 6--12 months.
\end{itemize}

\section{Key References}
\begin{itemize}
    \item American Academy of Ophthalmology (AAO). \textit{Preferred Practice Pattern: Refractive Errors \& Refractive Surgery} and \textit{Preferred Practice Pattern: Corneal Ectasia}.
    \item Rabinowitz YS. \textit{Videokeratographic Indices to Aid in Screening for Keratoconus}. J Refract Surg; and Gomes JAP, et al. \textit{Global Consensus on Keratoconus and Ectatic Diseases}. Cornea.
    \item American National Standards Institute (ANSI). \textit{ANSI Z80.23: Ophthalmic Instruments---Corneal Topography Systems}.
    \item StatPearls Publishing. \textit{Corneal Topography}.
    \item Cleveland Clinic. \textit{Corneal Topography}. Cleveland Clinic. Available at: \url{https://my.clevelandclinic.org/health/diagnostics/24819-corneal-topography}.
\end{itemize}

\section{Further Reading}
\begin{itemize}
  \item Wikipedia contributors. \textit{Corneal topography}. Wikipedia, The Free Encyclopedia. Available at: \url{https://en.wikipedia.org/wiki/Corneal_topography}.
\end{itemize}
